% Quick reference (practice author): the book's last page, input after the references (PLAN.md sections 1.2, 1.4).
% It repeats by design the five-row method choice T0 of section 1.5 and the two lessons of the edge (style decision
% 16). Unnumbered; listed in the contents. Refer to it with \hyperref[ch:quick-reference]{...} or \pageref.
\chapter*{Quick reference}\phantomsection\label{ch:quick-reference}
\addcontentsline{toc}{chapter}{Quick reference}\markboth{Quick reference}{}
\begin{fullwidth}
\small\setstretch{1.0}\setlength{\parskip}{4pt plus 1pt}
\setlength{\columnsep}{8mm}
\begin{multicols}{2}
\raggedright

\qrhead{Include.} \texttt{\#include <numerixx/deriv.hpp>}, in the namespace \code{nxx::deriv}. Header-only, C++23;
the CMake target is \code{numerixx::numerixx} (\cref{sec:api-synopsis}).

\qrhead{The call} (\cref{sec:api-derivative}).
\begin{itemize}\setlength{\itemsep}{0pt}\setlength{\parskip}{0pt}
  \item \code{derivative(fn, x)}: \code{Central<2>}, $f'$, default step;
  \item \code{derivative(fn, x, h)}: your own step, \code{h > 0};
  \item \code{derivative<M>(fn, x)}: the method \code{M};
  \item \code{derivative<M, N>(fn, x)}, with or without \code{h}: the derivative order \code{N}, from 1 to 4, or
    \code{order::second} and so on.
\end{itemize}
\code{fn} takes a \code{T}. \code{x} fixes \code{T}, which can be any \code{nxx::real} type.

\qrhead{The result.} A \code{std::expected<Estimate<T>, error>}. \code{d->value} is the estimate.
\code{*d->abs_error} is the error estimate of the three Ridders methods (empty for the stencils). \code{d.error()} is
the reason for a failure. \code{d.value()} is the accessor of \code{std::expected}, and it throws on a failure: test
first, then read \code{d->value}.

\qrhead{Two checks} (\cref{sec:using-result}). \emph{Did it work?} Test \code{if (d)}. \emph{Is it right?} A wrong
step scale, a noisy $f$ or a break near $x$ returns a wrong value without failing. If in doubt, repeat at $2h$ and
$h/2$. For Ridders, also require \code{r->abs_error && *r->abs_error <= atol + rtol * std::abs(r->value)}. Here
\code{rtol} $= 10^{-n}$ for $n$ significant figures, and \code{atol}, in the derivative's units, decides near
$f' = 0$. A stencil has no estimate to test. Inside an iteration, no check is needed.

\qrhead{Failures.} \code{non_finite} when \code{x} is not finite; when \code{fn} is infinite or NaN at a sample; on an
overflow; when the step is below the resolution of \code{x}; and when the step is not positive and finite. A float
derivative above the first warns at compile time, by every method.

\qrhead{Methods} (relative accuracy in double at the default step).
\begin{itemize}\setlength{\itemsep}{1pt}\setlength{\parskip}{0pt}
  \item \code{Central<P>}, P = 2, 4, \ldots, 10: $P + N - 1$ evaluations; an error of about $\epsmach^{P/(P+N)}$.
    With P = 2, about \num{\valHeadlineFirstDigits} digits for $f'$ and \num{\valHeadlineFourthDigits} for $f''''$.
  \item \code{Forward<P>}, \code{Backward<P>}, P = 1 to 6: $P + N$ evaluations, never below $x$ (\code{Forward})
    or above it (\code{Backward}); larger errors than \code{Central<P>}.
  \item \code{Ridders}: a median of \num{\valRiddersEvalsMedianFirst} evaluations for $f'$, and at most
    \num{\valRiddersMaxEvalsFirst}, \num{\valRiddersMaxEvalsSecond}, \num{\valRiddersMaxEvalsThird} and
    \num{\valRiddersMaxEvalsFourth} for $f'$ to $f''''$; about \num{\valHeadlineRiddersFirstDigits} digits for $f'$;
    an error estimate. Near a NaN, it returns a value down to
    $2^{-\num{\valRiddersEdgeLimitExponent}} \max(\abs{x}, 1)$ away (ten digits down to
    $2^{-\num{\valRiddersEdgeTenDigitsExponent}} \max(\abs{x}, 1)$), and then \code{non_finite}. Closer to a pole,
    it returns a wrong value, whose estimate flags it.
  \item \code{ForwardRidders}, \code{BackwardRidders}: Ridders on one side of $x$ only, for an edge, with an error
    estimate. A median of \num{\valForwardRiddersEvalsMedianFirst} evaluations for $f'$, at most
    \num{\valForwardRiddersMaxEvalsFirst}, \num{\valForwardRiddersMaxEvalsSecond},
    \num{\valForwardRiddersMaxEvalsThird} and \num{\valForwardRiddersMaxEvalsFourth}; a relative error of about
    \num{\valExpForwardRiddersFirstRelError} for $f'$.
\end{itemize}

\columnbreak
\qrhead{Choosing a method} (\cref{tab:method-summary}; the full table is \cref{tab:method-choice}).\par\smallskip
{\footnotesize\setlength{\tabcolsep}{3pt}
\begin{tabular}{@{}>{\raggedright}p{27mm}>{\raggedright}p{26mm}>{\raggedright\arraybackslash}p{25mm}@{}}
\toprule
Situation & Method and step & Cost and accuracy\\
\midrule
\tablerows{chapters/shared-method-summary}
\bottomrule
\end{tabular}}

\qrhead{Your own step.} \code{default_step<M, N>(x)} is the step that \code{derivative} uses. For
$\epsmach^{1/3} w$, write \code{h = nxx::math::root_epsilon<3, double>() * w;} (\code{Central<2>}, $f'$; in general
\code{root_epsilon<P + N, T>()}). For a noisy $f$, write
\code{h = std::cbrt(eta) * std::max(std::abs(x), 1.0);}.

\qrhead{No error estimate? Check it} (\cref{sec:using-check}). Repeat at $2h$ and $h/2$, and count the digits that
agree. Or compare with \code{Central<4>}, or with Ridders' method. Where $f' \approx 0$, judge the absolute error.
Check the sign against the physics.

\qrhead{The pitfalls that matter most} (\cref{sec:using-pitfalls}).
\begin{itemize}\setlength{\itemsep}{1pt}\setlength{\parskip}{0pt}
  \item \emph{The step's scale is in your units.} The default step assumes that $f$ varies on the scale
    $\max(\abs{x}, 1)$. At a large $\abs{x}$ with a narrow feature, or at $\abs{x} \ll 1$, pass $h$.
  \item \emph{Evaluable is not valid.} A fit sampled outside its range still returns numbers. Make it return NaN
    there. Then use \code{Forward<P>} or \code{Backward<P>} at the edge, or \code{ForwardRidders} or
    \code{BackwardRidders} for an error estimate.
  \item \emph{Evaluable is not smooth.} A kink or a range break within a step of $x$ gives a wrong value, without
    failing, from every method that samples across it. Ridders' estimate flags a jump; nothing flags a kink at $x$.
    Differentiate the branch you mean.
  \item \emph{\code{auto} and \code{.value()}.} \code{auto e = *d} is an \code{Estimate}, and
    \code{d.value()} throws on a failure. Write \code{d->value}.

\end{itemize}

\qrhead{Where to look.} Names: \cref{ch:api}; tables: \cref{ch:tables}; accuracy: \cref{tab:accuracy-matrix}; cost:
\cref{tab:cost}; guarantees: \cref{sec:guarantees-guarantees}.


\end{multicols}
\end{fullwidth}
